% 図2：ラベル伝播のグラフ．(a) 節点と辺と質量の流れ，(b) 再出現する未知の集合 Q とカット w(Q, \bar Q)，(c) 値の尺度はビュー・テスト集合で変わるが較正順位は同じ物差し．
% 全体の幅は約 17.3 cm（\textwidth = 17.7 cm），高さ約 4.8 cm．文字は \figfont（6.74pt）．
\providecommand{\figfont}{\fontsize{6.74pt}{8pt}\selectfont}
\begin{tikzpicture}[font=\figfont, >={Latex[length=1.3mm]}, x=1cm, y=1cm,
  sup/.style={draw=black, fill=black!75, minimum size=0.22cm, inner sep=0pt},
  idsq/.style={draw=black!55, fill=black!10, minimum size=0.2cm, inner sep=0pt},
  oodA/.style={draw=blue!60!black, fill=blue!35, minimum size=0.23cm, inner sep=0pt, regular polygon, regular polygon sides=3},
  calp/.style={draw=blue!70!black, fill=white, line width=0.6pt, minimum size=0.18cm, inner sep=0pt, circle},
  edge/.style={black!45, line width=0.45pt},
  flow/.style={->, orange!85!black, line width=0.7pt},
  lab/.style={inner sep=1pt, font=\figfont}]
% ================================================================ (a) graph and mass flow
\node[lab, anchor=west, font=\figfont\bfseries] at (0,4.6) {(a) 伝播のグラフ：支持画像の質量 $u^{+}$ が辺を伝わる};
\node[sup] (s1) at (0.7,3.85) {}; \node[sup] (s2) at (1.25,4.1) {}; \node[sup] (s3) at (1.7,3.7) {};
\node[lab, anchor=west] at (1.85,4.15) {支持画像（$y_j{=}1$）};
\node[calp] (c1) at (1.05,3.4) {}; \node[calp] (c2) at (2.1,3.2) {};
\node[idsq] (t1) at (0.5,3.1) {}; \node[idsq] (t2) at (1.45,2.95) {}; \node[idsq] (t3) at (2.45,3.85) {};
\node[idsq] (t4) at (2.85,2.15) {}; \node[calp] (c3) at (3.2,2.85) {};
\node[oodA] (q1) at (3.8,2.35) {}; \node[oodA] (q2) at (4.4,2.65) {}; \node[oodA] (q3) at (4.2,1.95) {}; \node[oodA] (q4) at (4.85,2.2) {};
\foreach \a/\b in {s1/s2, s2/s3, s1/t1, s1/c1, s2/c1, s3/t2, s3/c2, s3/t3, c1/t1, c1/t2, t2/c2, t3/c2, t2/t4, c2/c3, t4/c3, c3/q1, t4/q1, q1/q2, q1/q3, q2/q3, q2/q4, q3/q4, q1/q4}{\draw[edge] (\a) -- (\b);}
\draw[flow] (s1) to[bend right=12] (t1);
\draw[flow] (s3) to[bend left=12] (t3);
\draw[flow] (s3) to[bend right=10] (t2);
\draw[flow, orange!60!black, line width=0.5pt] (t2) to[bend right=8] (t4);
\draw[flow, orange!50!black, line width=0.35pt, densely dashed] (t4) -- (q1);
\node[lab, anchor=north, align=center] at (1.0,2.7) {ID：支持画像に近い\\$u^{+}$ 大 $\Rightarrow p^{+}$ 大};
\node[lab, anchor=north, align=center] at (4.35,1.72) {同じ未知クラス：互いに近い\\$u^{+}$ 小 $\Rightarrow p^{+}$ 小};
\node[lab, anchor=west] at (0.0,0.8) {\tikz{\node[sup]{};}\,支持 \tikz{\node[calp]{};}\,参照・較正 $\mathcal{R},\mathcal{C}$ \tikz{\node[idsq]{};}\,テスト（ID） \tikz{\node[oodA]{};}\,テスト（未知）};
\node[lab, anchor=west, orange!85!black] at (0.0,0.4) {$\rightarrow$ $u^{+}$ の流れ（\eqnref{lp}）．$u^{-}$ も同じ辺を流れる};
% ================================================================ (b) the cut
\begin{scope}[xshift=5.85cm]
\node[lab, anchor=west, font=\figfont\bfseries] at (0,4.6) {(b) 再出現する未知の集合 $Q$ とカット $w(Q,\bar{Q})$};
\node[oodA] (p1) at (1.6,3.2) {}; \node[oodA] (p2) at (2.25,3.55) {}; \node[oodA] (p3) at (2.15,2.8) {}; \node[oodA] (p4) at (2.8,3.15) {}; \node[oodA] (p5) at (2.75,2.5) {};
\node[idsq] (o1) at (0.6,3.8) {}; \node[idsq] (o2) at (0.55,2.6) {}; \node[idsq] (o3) at (3.9,3.8) {}; \node[sup] (o4) at (4.0,2.25) {};
\node[lab, anchor=west] at (4.15,2.25) {支持};
\foreach \a/\b in {p1/p2, p1/p3, p2/p3, p2/p4, p3/p4, p3/p5, p4/p5, p1/p5}{\draw[edge, line width=0.75pt] (\a) -- (\b);}
\draw[edge, densely dotted, line width=0.7pt] (o1) -- (p1); \draw[edge] (o2) -- (p1); \draw[edge, densely dotted, line width=0.7pt] (o3) -- (p4); \draw[edge] (o4) -- (p5);
\draw[red!70!black, densely dashed, line width=0.6pt] (2.2,3.1) ellipse (1.0cm and 0.8cm);
\node[lab, red!70!black, anchor=south] at (2.2,3.95) {$Q$（同じ未知クラス $r$ 枚）};
\node[lab, anchor=west, align=left] at (0.0,1.85) {$Q$ が自ら張る辺（実線）は $r$ とともに減る．外の節点が\\$Q$ を選ぶ辺（点線）は残る．};
\node[lab, anchor=west] at (0.0,1.2) {$\displaystyle\sum_{i\in Q}\sqrt{\delta_i}\,u_i\le\frac{\lambda}{1-\lambda}\cdot\frac{w(Q,\bar{Q})}{\min_{j\in Y}\sqrt{\delta_j}}$\quad（命題2）};
\node[lab, anchor=west, align=left] at (0.0,0.42) {$Q$ の各節点が外へ張る辺は $(k_g-r+1)_+$ 本．\\$w(Q,\bar{Q})$ はこれと点線の和で，再出現は\\前者だけを減らす（$Y{=}\hat{Q}$，$Q{=}$ID なら $u^{-}$ の上界）};
\end{scope}
% ================================================================ (c) scale changes, the calibrated rank does not
\begin{scope}[xshift=11.7cm]
\node[lab, anchor=west, font=\figfont\bfseries] at (0,4.6) {(c) $u$ の尺度はビューやテスト集合で変わる};
\node[lab, blue!70!black, anchor=west] at (0.0,4.15) {順位で読めば閾値（$\mathcal{R}$ の下位 $\varepsilon$ 点）も一緒に動く};
\node[lab, anchor=east] at (1.05,3.5) {ビュー $v_1$};
\node[lab, anchor=east] at (1.05,2.4) {ビュー $v_2$};
\draw[->, black!45] (1.15,3.3) -- (5.0,3.3); \node[lab, anchor=north east, black!65] at (5.05,3.28) {$u^{+}$（右ほど ID らしい）};
\draw[->, black!45] (1.15,2.2) -- (5.0,2.2);
\foreach \x in {2.45,2.9,3.25,3.75,4.1,4.45} {\node[calp] at (\x,3.5) {};}
\node[idsq] at (3.55,3.5) {}; \node[oodA] at (1.6,3.51) {};
\foreach \x in {1.85,2.1,2.3,2.6,2.8,3.0} {\node[calp] at (\x,2.4) {};}
\node[idsq] at (2.48,2.4) {}; \node[oodA] at (1.45,2.41) {};
\draw[blue!70!black, line width=0.8pt] (3.1,3.8) -- (3.1,3.25); \draw[blue!70!black, line width=0.8pt] (2.2,2.7) -- (2.2,2.15);
\node[lab, blue!70!black, anchor=west] at (3.15,3.8) {$\mathcal{R}$ の下位 $\varepsilon$ 点};
\draw[red!70!black, line width=0.8pt] (2.65,3.05) -- (2.65,2.05);
\node[lab, red!70!black, anchor=west, align=left] at (1.15,1.7) {$u^{+}$ の固定閾値（赤線）：$v_2$ では ID が\\下回る．ビューの積も取れない};
\node[lab, anchor=west, align=left] at (0.0,0.6) {命題1：$\bom^{\!\top}\mathbf{u}^{+}=\sum_{j\in\mathcal{S}}\sqrt{\delta_j}$ は支持画像の次数\\だけで決まる（節点 $N$ が増えるほど $u^{+}$ は縮む）};
\end{scope}
\end{tikzpicture}
